% Part: many-valued-logic
% Chapter: three-valued-logics
% Section: introduction

\documentclass[../../../include/open-logic-section]{subfiles}

\begin{document}

\olfileid{mvl}{thr}{int}

\olsection{مقدمة}

إذا زدنا قيمة ثالثة هي~$\Undef$ على~$\True$ و$\False$، نشأ منطق
ثلاثي القيم. وليس انحصار الزيادة في قيمة واحدة موجبًا لانحصار دوال
الصدق؛ بل يمكن تعريف~$\lnot$ و$\land$ و$\lor$ و$\lif$ على وجوه
كثيرة جدًا. ولكل منطق أن يختار أي تركيب من تلك الدوال، وأن يجعل
$\True$ وحدها مميزة أو يجعل~$\True$ و$\Undef$ معًا مميزتين.

وسنعرض فيما يأتي جماعة من أشهر المنطقات الثلاثية القيم، مبينين ما
دعا إلى كل منها وبعض ما يختص به.

\end{document}
